\section{Introduction}

EEG foundation models promise reusable representations across heterogeneous
montages, recording systems, cohorts, and downstream tasks. REVE was designed
for this setting through large-scale masked pretraining and position encodings
that accommodate varying electrode layouts \citep{elouahidi2025reve}. At the
same time, a good downstream score does not by itself identify what a frozen
representation makes readily accessible: a sufficiently flexible head may learn
dataset-specific structure that a linear readout cannot exploit.

Age is a useful stress test. It is continuously valued, strongly reflected in
EEG physiology, and studied as a marker of maturation and aging
\citep{sun2019brainage,vandenbosch2019age}. Yet age decoding is also vulnerable
to acquisition and cohort shift. A result that improves on an internal split
may fail when hardware, recruitment, and recording details change.

We therefore ask a deliberately narrow question: when the REVE encoder,
development data, preprocessing, optimization seeds, and checkpoint rule are
held fixed, do alternative or richer representation heads improve reliably over
a mean-pooled linear probe? The study separates three evidence classes. The
primary evidence is a prospective frozen-probe comparison developed on HBN and
evaluated once on a sealed MIPDB cohort. Official NeuralBench-style end-to-end
experiments provide reproduction context, while previously inspected HBN test
results are retrospective only. MIPDB is not an official NeuralBench score.

We additionally examine whether the paired head comparison changes with the
amount of development data. A secondary capacity--data extension trains the
mean-linear baseline and two richer heads on nested HBN subsets of 200, 400,
and 800 subjects, using ten fixed optimization seeds and the same sealed MIPDB
evaluation. This extension is explicitly exploratory: it probes whether more
development data makes a richer head more reliable, but it does not expand the
primary confirmatory claim or represent a full-dataset scaling law.

Our contribution is methodological as much as numerical. We combine a matched
four-head design, exact repeated-seed inventories, subject- and seed-aware
uncertainty, a paired seed-level sign-flip test with multiplicity correction,
and a conjunctive stability criterion. We then use the secondary extension to
separate absolute performance from changes in the paired candidate--baseline
difference across training regimes. The outcome illustrates why average gains
and small adjusted p-values should not be interpreted as stable representation
improvements when external uncertainty remains broad.
